\appendix
\section*{Appendix Q: Holonomy Trace of NDH–8B–StandardGalactic Cross-Pollination}

\subsection*{Q.0 Overview}

This appendix defines the holonomy trace framework used to analyze
cross-pollination between NDH primitives, 8B resonance algebra,
StandardGalactic conceptual surfaces, and Omnibus ecosystem reflections.
The analysis is performed inside the Fifty-Dimensional Neutral Manifold
$M^{50}$ with geometric operators decomposed through the Unified Altitude
Algebra and validated by the Series V firewall.

The purpose is to:
\begin{itemize}
    \item extract holonomy signatures of conceptual primitives,
    \item identify curvature-driven resonance paths,
    \item detect parallel-transport distortions,
    \item characterize conceptual twist vectors,
    \item identify orientation-surface formation (Direction Manifold v2.0),
    \item provide holonomy operators for Julia runtime analysis.
\end{itemize}

\subsection*{Q.1 Holonomy Embedding in the 50D Neutral Manifold}

Let:
\[
\iota_{NDH},\; \iota_{8B},\; \iota_{SG} : X \to M^{50}
\]
be altitude-safe embeddings of conceptual primitives into the neutral
containment manifold.

Each embedded element $x$ admits a holonomy representation:
\[
\mathrm{Hol}(x) = \exp\left(\oint_{\gamma_x} \omega\right)
\]
where $\omega$ is the connection 1-form associated with the unified
algebra and $\gamma_x$ is a conceptual loop in $M^{50}$.

\subsection*{Q.2 Holonomy Operators}

Define holonomy operator:
\[
H : M^{50} \to GL(50)
\]
such that:
\[
H(x) = \mathrm{Hol}(x).
\]

For each ecosystem:
\[
H_{NDH},\quad H_{8B},\quad H_{SG}.
\]

Cross-pollination is detected when:
\[
H_{NDH}(x) \approx H_{8B}(y)
\quad\text{or}\quad
H_{NDH}(x) \approx H_{SG}(z).
\]

\subsection*{Q.3 Curvature Envelopes}

Define curvature envelope:
\[
K(x) = d\omega + \omega \wedge \omega.
\]

Curvature alignment indicates:
\begin{itemize}
    \item conceptual resonance,
    \item orientation-surface formation,
    \item primitive migration,
    \item ecosystem twist absorption.
\end{itemize}

Cross-pollination curvature is detected when:
\[
K_{NDH}(x) \approx K_{8B}(y)
\quad\text{or}\quad
K_{NDH}(x) \approx K_{SG}(z).
\]

\subsection*{Q.4 Parallel Transport Distortion}

Define parallel transport operator:
\[
T_{\gamma}(x) = \mathcal{P}\exp\left(\int_{\gamma} \omega\right).
\]

Distortion is measured by:
\[
\Delta_{\gamma}(x,y)
=
\|T_{\gamma}(x) - T_{\gamma}(y)\|.
\]

Cross-pollination distortion occurs when:
\[
\Delta_{\gamma}(x,y) < \epsilon.
\]

\subsection*{Q.5 Conceptual Twist Vectors}

Define twist vector:
\[
\tau(x) = \mathrm{Hol}(x) - I.
\]

Twist alignment indicates:
\begin{itemize}
    \item shared conceptual primitives,
    \item resonance absorption,
    \item atlas-level synthesis,
    \item softmax-driven conceptual blending.
\end{itemize}

\subsection*{Q.6 Orientation Surface Formation}

Let $\mathcal{D}_2$ be the Direction Manifold v2.0.

Define orientation surface formation condition:
\[
\mathrm{Hol}(x) \in \mathrm{Span}(\mathcal{D}_2)
\]

Cross-pollination orientation surfaces occur when:
\[
\mathrm{Hol}_{NDH}(x),\;
\mathrm{Hol}_{8B}(y),\;
\mathrm{Hol}_{SG}(z)
\in
\mathrm{Span}(\mathcal{D}_2).
\]

\subsection*{Q.7 Holonomy Projection Towers}

Define holonomy projection tower:
\[
P_{\mathrm{Hol}}(x)
=
(P_1(x), P_2(x), \dots, P_{50}(x)).
\]

Projection alignment indicates:
\begin{itemize}
    \item curvature compatibility,
    \item holonomy resonance,
    \item conceptual twist coherence.
\end{itemize}

\subsection*{Q.8 Firewall Enforcement}

All holonomy comparisons are wrapped in:
\[
\mathrm{Firewall}_{\mathrm{unified}}(a,b,T)
\]
ensuring:
\begin{itemize}
    \item purity,
    \item non-generativity,
    \item boundedness,
    \item altitude legality.
\end{itemize}

\subsection*{Q.9 Output for Julia Runtime}

Julia receives:
\begin{itemize}
    \item holonomy operators $H$,
    \item curvature envelopes $K$,
    \item parallel transport operators $T_{\gamma}$,
    \item twist vectors $\tau$,
    \item holonomy projection towers $P_{\mathrm{Hol}}$.
\end{itemize}

These form the geometric analysis layer of the runtime experiment.

This completes the holonomy trace appendix.
